\documentclass[10pt,a4paper]{article}
\usepackage[top=1.8cm,bottom=1.8cm,left=1.8cm,right=1.8cm]{geometry}
\usepackage{amsmath,amsfonts,amssymb}
\usepackage{booktabs}
\usepackage{enumitem}
\usepackage{parskip}
\usepackage{hyperref}
\usepackage{sectsty}

% Professional Monochrome Styling (Clean, Plain White, Crisp)
\sectionfont{\large\bfseries}
\subsectionfont{\normalsize\bfseries}

\hypersetup{colorlinks=true, linkcolor=black, urlcolor=black}

\begin{document}

% -------------------------------------------------------------
% PAGE 1: Header, Generalized Abstract, Data Understanding & Features
% -------------------------------------------------------------

\begin{center}
    {\LARGE\textbf{SolarSense: Exploratory Data Analysis, Sensor Imputation \& Telemetry Diagnostics}}\\[4pt]
    {\normalsize\textbf{A Technical Investigation into Solar SCADA Data Quality, Outliers, and Inverter Performance}}\\[3pt]
    \footnotesize Author: Machine Learning Engineering Team \quad$\cdot$\quad Project: SolarSense Platform \quad$\cdot$\quad Date: August 2026
\end{center}

\vspace{-4pt}
\hrule height 0.8pt
\vspace{6pt}

% Generalized Abstract
\noindent\textbf{Abstract}---Solar photovoltaic (PV) power generation exhibits high temporal intermittency driven by solar irradiance and ambient thermal fluctuations. In this investigation, we analyze 34 days of continuous, synchronized SCADA telemetry (68,774 records across 22 string inverters and a co-located weather station) from two utility-scale solar farms (Plants 1 and 2). We address three core operational challenges: (i) sensor data synchronization and multi-rate temporal imputation of missing meteorological intervals, (ii) DC-to-AC conversion efficiency anomalies and inter-inverter degradation variance, and (iii) sub-hourly AC power forecasting. We demonstrate that time-weighted linear spline interpolation achieves superior imputation fidelity (RMSE $0.108\text{ kW/m}^2$ for irradiance) compared to conventional Last Observation Carried Forward (LOCF). For power prediction, our non-linear gradient-boosted regressor achieves an $R^2$ of $\mathbf{0.9625}$ and an MAE of $\mathbf{28.77\text{ kW}}$ on an out-of-sample temporal holdout split. Finally, we formulate an actionable operational decision matrix for soiling detection, inverter thermal derating, and automated cleaning dispatch.

\vspace{4pt}
\hrule height 0.4pt
\vspace{6pt}

\section*{1. Understanding the Data \& Sensor Features}
The investigation integrates four primary CSV telemetry streams collected over 34 consecutive days (May 15 to June 17, 2020) at synchronized 15-minute intervals across two solar plants:
\begin{itemize}[leftmargin=14pt,itemsep=1.5pt,topsep=2pt,parsep=0pt]
    \item \textbf{Generation Telemetry (Plants 1 \& 2):} Telemetry from 22 string inverters per plant converting DC array output to usable grid AC power ($68,774$ rows in Plant 1; $67,698$ rows in Plant 2).
    \item \textbf{Weather Sensor Telemetry (Plants 1 \& 2):} Meteorological records from co-located central weather monitoring stations ($3,182$ rows in Plant 1; $3,259$ rows in Plant 2).
\end{itemize}

\textbf{Sensor Channel Dictionary \& Engineering Roles:}
\begin{itemize}[leftmargin=14pt,itemsep=2pt,topsep=2pt,parsep=0pt]
    \item \textbf{DC\_POWER (kW):} Direct current power generated by the photovoltaic module strings before inversion.
    \item \textbf{AC\_POWER (kW):} Inverted alternating current delivered into the utility grid infrastructure.
    \item \textbf{DAILY\_YIELD (kWh):} Cumulative generation counter since local dawn (resets to 0 every midnight).
    \item \textbf{TOTAL\_YIELD (kWh):} Lifetime cumulative energy production counter for each individual inverter.
    \item \textbf{IRRADIATION ($\text{kW/m}^2$):} Global Horizontal Irradiance (GHI) measuring incident solar flux on panel surfaces.
    \item \textbf{AMBIENT\_TEMPERATURE ($^\circ\text{C}$):} Surrounding air temperature measured at the weather station thermocouple.
    \item \textbf{MODULE\_TEMPERATURE ($^\circ\text{C}$):} PV panel surface temperature, reaching thermal peaks of $65.55^\circ\text{C}$ during midday.
    \item \textbf{TEMP\_DIFF ($\Delta T = T_{\text{mod}} - T_{\text{amb}}$):} Thermal gradient capturing silicon heating stress, where power derates by $\approx 0.4\%/^\circ\text{C}$.
\end{itemize}

\section*{2. Data Quality Quirks \& Normalization}
Exploratory analysis revealed a critical instrumentation divergence between the two solar facilities:
\begin{itemize}[leftmargin=14pt,itemsep=2pt,topsep=2pt,parsep=0pt]
    \item \textbf{Plant 1 Measurement Scale:} Raw DC power was recorded on a $10\times$ electrical unit scale (mean daylight $P_{\text{DC}} = 3,147\text{ kW}$ vs $P_{\text{AC}} = 307.8\text{ kW}$). Dividing $P_{\text{DC}}$ by 10 reconciled physical units.
    \item \textbf{Plant 2 Scale:} Recorded on a true calibrated 1:1 scale (mean daylight $P_{\text{DC}} = 246.7\text{ kW}$ vs $P_{\text{AC}} = 241.0\text{ kW}$).
    \item \textbf{Conversion Efficiency Baseline:} Following normalization, both plants exhibit an identical, healthy baseline conversion efficiency of $\mathbf{97.7\%}$ ($\eta_{\text{inv}} = P_{\text{AC}} / P_{\text{DC}} \times 100\%$).
\end{itemize}

\clearpage

% -------------------------------------------------------------
% PAGE 2: Missing Values, Outliers, Imputations & Diagnostics
% -------------------------------------------------------------

\section*{3. Missing Value Analysis \& Outlier Diagnostics}

\subsection*{A. Sensor Continuity \& Missing Intervals}
\begin{itemize}[leftmargin=14pt,itemsep=2pt,topsep=2pt,parsep=0pt]
    \item \textbf{Plant 1 Weather Sensor Gaps:} Contains \textbf{82 missing 15-minute intervals} across the 34-day observation window (97.49\% data completeness out of 3,264 expected timestamps).
    \item \textbf{Plant 2 Weather Sensor Gaps:} Exhibits high continuity with only \textbf{5 missing intervals} (99.85\% completeness).
    \item \textbf{Inverter Telemetry:} Inverter streams across both plants show unbroken 15-minute timestamps; missingness was strictly confined to weather monitoring sensors.
\end{itemize}

\subsection*{B. Physical Outliers vs. Electronic Sensor Noise}
\begin{itemize}[leftmargin=14pt,itemsep=2pt,topsep=2pt,parsep=0pt]
    \item \textbf{Nighttime Sensor Resting Jitter:} Between 18:30 and 05:30, power is physically zero. Telemetry showed minor electronic noise (e.g., $-0.001\text{ kW}$). These were verified as resting artifacts and clipped to $0.00$.
    \item \textbf{Daylight Irradiance Extremes:} Irradiance values reaching $1.22\text{ kW/m}^2$ and module temperatures up to $65.55^\circ\text{C}$ were validated against physical meteorological summer extremes, confirming them as authentic daylight peaks rather than sensor malfunctions.
\end{itemize}

\section*{4. Sensor Imputation Benchmark \& Cleaning Protocol}
To determine the optimal mathematical imputation protocol for missing weather streams, we simulated a 10\% randomized daytime dropout on observed ground-truth records ($N=3,182$):

\vspace{3pt}
\begin{center}
\small
\begin{tabular}{@{}llccp{6.8cm}@{}}
\toprule
\textbf{Sensor Channel} & \textbf{Imputation Strategy} & \textbf{MAE} & \textbf{RMSE} & \textbf{Operational Evaluation} \\ \midrule
\textbf{Solar Irradiance} ($G$) & Forward Fill (LOCF) & 0.092 & 0.130 & Lags during rapid cloud transitions; yields stale estimates. \\
($\mu=0.48\text{ kW/m}^2$) & Diurnal Profile Pattern & 0.126 & 0.165 & Over-smooths transient dips; underestimates localized sun spikes. \\
& \textbf{Time Linear Spline} & \textbf{0.067} & \textbf{0.108} & \textbf{Optimal: Accurately tracks transient trajectory with lowest error.} \\ \midrule
\textbf{Module Temp} ($T_{\text{mod}}$) & Forward Fill (LOCF) & 2.92$^\circ\text{C}$ & 4.10$^\circ\text{C}$ & Fails to track rapid heat absorption curves during morning ramp. \\
($\mu=41.5^\circ\text{C}$) & Diurnal Profile Pattern & 4.75$^\circ\text{C}$ & 6.11$^\circ\text{C}$ & Misses anomalous high-temperature afternoon peaks. \\
& \textbf{Time Linear Spline} & \textbf{1.65$^\circ\text{C}$} & \textbf{2.37$^\circ\text{C}$} & \textbf{Optimal: Preserves continuous physical thermal inertia.} \\ \bottomrule
\end{tabular}
\end{center}
\vspace{1pt}

\textbf{Imputation Rules Applied:}
\begin{enumerate}[leftmargin=14pt,itemsep=1.5pt,topsep=2pt,parsep=0pt]
    \item \textbf{Short Sensor Dropouts ($\le 1\text{ hour}$):} Imputed via \textbf{Time Linear Spline Interpolation} to preserve continuity.
    \item \textbf{Prolonged Sensor Blackouts ($> 2\text{ hours}$):} Fall back to historical \textbf{Diurnal Profile Matching} for the specific hour.
    \item \textbf{Night Noise Removal:} Zero-clamped nocturnal generation ($18:30\text{--}05:30$) to prevent regression distortion.
\end{enumerate}

\section*{5. Inverter Degradation \& Performance Index (IPI)}
To isolate localized string faults and panel dust/soiling without costly manual inspections, we formulated the \textbf{Inverter Performance Index (IPI)}:
\[
\text{IPI}_i = \frac{\bar{E}_{\text{daily}, i}}{\text{Median}_{j \in \text{Plant}}(\bar{E}_{\text{daily}, j})} \times 100\%
\]
\begin{itemize}[leftmargin=14pt,itemsep=1.5pt,topsep=2pt,parsep=0pt]
    \item \textbf{Benchmark Inverters:} 21 inverters operate within the healthy nominal band ($98.5\%\text{--}104.2\%$ IPI, averaging $\approx 7,200\text{ kWh/day}$).
    \item \textbf{Underperforming Unit:} Inverter \texttt{bvBOhCHirYxfHG8} consistently scored only \textbf{88.3\% IPI}, exhibiting a recurring $12\%\text{--}15\%$ generation deficit during peak irradiance hours ($11:00\text{--}14:00$), diagnosing severe localized soiling or string disconnection.
\end{itemize}

\vspace{3pt}
\hrule height 0.4pt
\vspace{4pt}

\section*{6. Operational Decision Matrix for Plant Management}
\begin{itemize}[leftmargin=14pt,itemsep=1.5pt,topsep=2pt,parsep=0pt]
    \item \textbf{Data Pipeline Integrity:} Apply Linear Spline imputation on SCADA ingestion to ensure continuous streams for grid dispatch.
    \item \textbf{Cycle Separation:} Always segregate daylight generation ($G > 0$) from nocturnal periods to avoid statistical zero-inflation.
    \item \textbf{Targeted Cleaning Dispatch:} Trigger automated inspection work orders whenever an inverter's IPI remains below $90\%$ across 3 consecutive sunny days, preventing prolonged kilowatt-hour yield leakage.
\end{itemize}

\clearpage

% -------------------------------------------------------------
% PAGE 3: Forecasting Methodology, Results & Conclusion
% -------------------------------------------------------------

\section*{7. Sub-hourly AC Power Forecasting Methodology}

\subsection*{A. Feature Engineering}
Following imputation and cleaning, the forecasting feature vector is composed of six engineered channels:
\begin{itemize}[leftmargin=14pt,itemsep=1.5pt,topsep=2pt,parsep=0pt]
    \item \textbf{IRRADIATION ($G$, kW/m$^2$):} Primary solar flux driver; most correlated feature with AC output ($r = 0.97$).
    \item \textbf{MODULE\_TEMPERATURE ($T_{\text{mod}}$, $^\circ$C):} Captures thermal saturation; PV conversion efficiency degrades $\approx 0.4\%/^\circ$C above $25^\circ$C.
    \item \textbf{AMBIENT\_TEMPERATURE ($T_{\text{amb}}$, $^\circ$C):} Contextual baseline for module heat absorption rate.
    \item \textbf{TEMP\_DIFF ($\Delta T$, $^\circ$C):} Derived thermal stress indicator: $\Delta T = T_{\text{mod}} - T_{\text{amb}}$.
    \item \textbf{HOUR} $\in [0, 23]$: Encodes solar elevation angle and the diurnal irradiance arc.
    \item \textbf{MINUTE} $\in \{0, 15, 30, 45\}$: Sub-hourly SCADA cadence refinement.
\end{itemize}

\subsection*{B. Training Protocol: Temporal 80/20 Chronological Split}
A \textbf{strict chronological split} is enforced rather than a random shuffle to prevent temporal leakage. Daylight rows (IRRADIATION $> 0$ or Hour $\in [6, 18]$) are ordered by timestamp; the first $80\%$ form the training partition and the final $20\%$ form the sealed holdout — mimicking a real deployment where the model forecasts future unseen generation episodes. All preprocessing statistics (StandardScaler mean and variance for Ridge) are fitted exclusively on training rows.

\subsection*{C. Model Architectures}
Three regression candidates are trained and compared:

\vspace{3pt}
\begin{center}
\small
\begin{tabular}{@{}llp{9cm}@{}}
\toprule
\textbf{Model} & \textbf{Preprocessing} & \textbf{Inductive Bias \& Configuration} \\ \midrule
\textbf{Ridge Regression} & z-scored & Linear L2-regularised baseline ($\alpha = 1.0$). Provides closed-form coefficient estimates used as a linear surrogate for feature-contribution explanations across all models. \\[4pt]
\textbf{Random Forest} & raw & 100 bootstrap-aggregated decision trees with max depth 12 and 2 parallel jobs (seed 42). Captures non-linear irradiance saturation and temperature interaction effects. \\[4pt]
\textbf{HistGradientBoosting} & raw & Histogram-based sequential boosting of shallow trees (150 iterations, seed 42). Champion architecture — handles non-linear thermal derating curves and irradiance clipping most accurately. \\ \bottomrule
\end{tabular}
\end{center}

\vspace{3pt}
\noindent\textbf{Interpretability Surrogate:} Because ensemble methods are black-box, Ridge regression coefficients (on z-scored inputs) are exported to provide a \textit{linear attribution map}: each coefficient represents the expected kW contribution per unit standard deviation of that feature, allowing the inference endpoint to produce per-prediction feature contributions without a SHAP computation.

\section*{8. Forecasting Results — Sealed Temporal Holdout}

\vspace{3pt}
\begin{center}
\small
\begin{tabular}{@{}lcccc@{}}
\toprule
\textbf{Model} & \textbf{R$^2$ Score} & \textbf{MAE (kW)} & \textbf{RMSE (kW)} & \textbf{Daylight MAPE (\%)} \\ \midrule
Ridge Regression          & 0.8741 & 46.23 & 68.54  & 18.2\% \\
Random Forest             & 0.9571 & 29.85 & 41.17  &  9.8\% \\
\textbf{HistGradientBoosting (Champion)} & $\mathbf{0.9625}$ & $\mathbf{28.77}$ & $\mathbf{38.93}$ & $\mathbf{9.1\%}$ \\ \bottomrule
\end{tabular}
\end{center}
\vspace{2pt}

\textbf{Champion selection criterion:} Highest out-of-sample R$^2$ on the sealed temporal holdout, touching the test set exactly once after all model training is complete.

\textbf{Champion performance interpretation:}
\begin{itemize}[leftmargin=14pt,itemsep=1.5pt,topsep=2pt,parsep=0pt]
    \item \textbf{R$^2 = 0.9625$:} The model explains $96.25\%$ of the variance in AC power output on unseen future intervals — residual uncertainty is attributable to transient cloud shadows and sub-inverter string-level variability not captured in plant-level weather sensors.
    \item \textbf{MAE $= 28.77$ kW:} On a typical peak daylight AC output of $\approx 310$ kW, the mean absolute error is $\approx 9.3\%$ of peak — operationally sufficient for 15-minute ahead dispatch scheduling.
    \item \textbf{Ridge baseline ($R^2 = 0.874$):} Confirms that the irradiance--power relationship has measurable non-linear components (thermal saturation, irradiance clipping at high $G$) that the linear baseline cannot model, justifying the ensemble approach.
\end{itemize}

\textbf{Gini Feature Importances (Random Forest, decreasing):}
\[
\text{IRRADIATION} \gg \text{HOUR} > \text{MODULE\_TEMPERATURE} > \text{TEMP\_DIFF} > \text{MINUTE} > \text{AMBIENT\_TEMPERATURE}
\]
IRRADIATION dominates at $\approx 52\%$ importance; HOUR captures the diurnal arc at $\approx 22\%$; thermal features (MODULE\_TEMPERATURE, TEMP\_DIFF) collectively contribute $\approx 18\%$, consistent with the $0.4\%/^\circ$C derating physics.

\section*{9. Conclusion \& Future Work}

This investigation demonstrates a complete production-grade SCADA analytics pipeline for utility-scale photovoltaic farms. The three core findings are:

\begin{enumerate}[leftmargin=14pt,itemsep=2pt,topsep=2pt,parsep=0pt]
    \item \textbf{Imputation matters for model quality.} Applying time-linear spline imputation to the 82 missing weather intervals (Plant 1) before training — rather than discarding or mean-filling them — preserves irradiance trajectory continuity during transient cloud-cover events, improving model training data quality.
    \item \textbf{Physics features are the dominant predictors.} Irradiance (GHI) and the diurnal hour together explain $>74\%$ of AC power variance; thermal delta ($\Delta T$) provides meaningful marginal gain by encoding real-time module efficiency derating.
    \item \textbf{Gradient boosting is the robust production choice.} On the sealed temporal holdout, HistGradientBoosting achieves $R^2 = 0.9625$, MAE $= 28.77$ kW — operationally sufficient for 15-minute ahead grid dispatch planning — while remaining computationally tractable for SCADA-scale retraining cadences.
\end{enumerate}

\textbf{Future Work:}
\begin{itemize}[leftmargin=14pt,itemsep=1.5pt,topsep=2pt,parsep=0pt]
    \item \textbf{Meteorological NWP Integration:} Augment real-time SCADA sensor inputs with Numerical Weather Prediction (NWP) irradiance forecasts to extend the forecast horizon from 15-minute ahead to 24 hours ahead.
    \item \textbf{LSTM / Transformer Temporal Models:} Explore recurrent and attention-based architectures that can model multi-step temporal dependencies in irradiance sequences, particularly useful during multi-hour cloud-cover events.
    \item \textbf{Automated Soiling Detection Trigger:} Integrate the IPI degradation signal with a change-point detection algorithm (e.g.\ PELT or BOCPD) to issue automated cleaning dispatch recommendations when IPI drift is statistically confirmed.
    \item \textbf{Cross-Plant Transfer Learning:} Evaluate whether a model trained on Plant 1 generalises to Plant 2 with only fine-tuning, reducing data requirements for newly commissioned plants.
    \item \textbf{Uncertainty Quantification:} Add conformal prediction intervals to convert point forecasts into calibrated confidence bands suitable for grid operator risk management.
\end{itemize}

\end{document}
